% !TEX program = pdflatex
% Standalone output assembled by external_calibration_review_evidence/build_review.py.
\documentclass[10pt,a4paper,landscape]{article}
\usepackage[T1]{fontenc}
\usepackage[utf8]{inputenc}
\usepackage{lmodern,amsmath,amssymb}
\usepackage[margin=14mm,headheight=14pt,headsep=6mm,footskip=8mm]{geometry}
\usepackage{array,longtable,booktabs,ragged2e,multicol}
\usepackage[table]{xcolor}
\usepackage{microtype,fancyhdr}
\usepackage[hidelinks,bookmarksnumbered]{hyperref}
\definecolor{heading}{RGB}{28,53,73}
\definecolor{tablehead}{RGB}{231,238,243}
\pagestyle{fancy}\fancyhf{}
\fancyhead[L]{\small\sffamily\color{heading}Two-country HKT: external-calibration literature review}
\fancyhead[R]{\small\sffamily 27 September 2026}
\fancyfoot[C]{\thepage}
\renewcommand{\headrulewidth}{0.3pt}
\setlength{\parindent}{0pt}\setlength{\parskip}{5pt}
\setlength{\tabcolsep}{5pt}\setlength{\LTpre}{5pt}\setlength{\LTpost}{7pt}
\renewcommand{\arraystretch}{1.1}\setlength{\emergencystretch}{3em}
\newcolumntype{P}[1]{>{\raggedright\arraybackslash}p{#1}}
\newcommand{\src}[1]{\hyperlink{ref:#1}{\textsf{[#1]}}}
\newcommand{\code}[1]{{\small\texttt{#1}}}
\newenvironment{evidence}{%
\fontsize{9.4}{10.9}\selectfont\setlength{\parskip}{2pt}
\begin{longtable}{@{}P{\dimexpr .22\textwidth-7.5pt\relax}P{\dimexpr .24\textwidth-7.5pt\relax}P{\dimexpr .25\textwidth-7.5pt\relax}P{\dimexpr .29\textwidth-7.5pt\relax}@{}}
\toprule\rowcolor{tablehead}\textbf{Source and original object} & \textbf{Reported values / evidence type} & \textbf{Method, population and frequency} & \textbf{V6 transformation and transferability}\\\midrule
\endfirsthead
\toprule\rowcolor{tablehead}\textbf{Source and original object} & \textbf{Reported values / evidence type} & \textbf{Method, population and frequency} & \textbf{V6 transformation and transferability}\\\midrule
\endhead
\midrule\multicolumn{4}{r}{\footnotesize\itshape Continued on next page}\\\endfoot
\bottomrule\endlastfoot
}{\end{longtable}}
\newenvironment{recommendations}{%
\fontsize{9.4}{10.9}\selectfont\setlength{\parskip}{2pt}
\begin{longtable}{@{}P{\dimexpr .14\textwidth-8pt\relax}P{\dimexpr .25\textwidth-8pt\relax}P{\dimexpr .32\textwidth-8pt\relax}P{\dimexpr .15\textwidth-8pt\relax}P{\dimexpr .14\textwidth-8pt\relax}@{}}
\toprule\rowcolor{tablehead}\textbf{Parameter} & \textbf{Proposed baseline candidate} & \textbf{Robustness range / scenario meaning} & \textbf{Mapping confidence} & \textbf{Best support}\\\midrule
\endfirsthead
\toprule\rowcolor{tablehead}\textbf{Parameter} & \textbf{Proposed baseline candidate} & \textbf{Robustness range / scenario meaning} & \textbf{Mapping confidence} & \textbf{Best support}\\\midrule
\endhead
\midrule\multicolumn{5}{r}{\footnotesize\itshape Continued on next page}\\\endfoot
\bottomrule\endlastfoot
}{\end{longtable}}
\hypersetup{pdftitle={V6 external calibration: evidence, mappings and candidate ranges},pdfauthor={}}
\begin{document}
{\LARGE\sffamily\bfseries\color{heading}External calibration: evidence and candidate ranges}\par
{\large\sffamily V6 preference, production and technology parameters}\par

\textbf{Assessment.} The literature supplies useful conditional candidates for risk curvature,
markups and substitution elasticities. It does not supply an empirically defensible single
vector for all V6 parameters. Discounting requires a generation-length bridge; technology
levels require joint normalization and production-data inversion; innovation coefficients
and the absorbing technology regime require additional measurement or explicit scenarios.
The recommendations at the end distinguish numeric candidates from
quantities deliberately left data-dependent or unassigned.

\textbf{Authority and coverage.}
\code{V6\_parameter\_calibration\_groups.tex} defines the parameters reviewed here.
Coverage is \textbf{11 families / 19 scalar parameters}: M01, M02, M07, M09--M16.
M08 labor endowments are data inputs and receive no conventional-value search.
Group C, M03--M06, is reserved for internal international-finance calibration.
M17's $\nu_b=\xi_W=0$ are imposed restrictions, not free external parameters.
The current manuscript supplies supporting equations; no model definition or solver has been changed.

\begin{center}\small
\begin{tabular}{@{}P{.20\textwidth}P{.73\textwidth}@{}}
\toprule\rowcolor{tablehead}\textbf{Evidence label} & \textbf{Interpretation in this review}\\\midrule
Estimated & Original econometric or experimental evidence, conditional on its sample and model. Standard errors, confidence intervals and population percentiles are identified separately.\\[3pt]
Calibrated / illustrative & A value imposed by an author, including HKT's theoretical example; it is not an independent estimate.\\[3pt]
Synthesis & A meta-analysis of original studies; useful for bias and heterogeneity, without treating pooled estimates as a consensus.\\[3pt]
Proposed scenario & Our mapping or robustness choice; its endpoints are not a literature confidence interval. Unassigned means no portable numerical benchmark was verified.\\
\bottomrule
\end{tabular}
\end{center}

\textbf{Search and verification.}
Consensus searches and fetched records, Scite searches/full text, and an
\href{https://app.undermind.ai/projects/6513b93c-0aaf-4fd9-9554-e57b8db4a0cf?path=/External\%20calibration\%20of\%20V6\%20preferences\%20and\%20production}{Undermind deep search}
were used for discovery, followed by original papers and dated author versions for numerical checks.
The Undermind search completed with 175 ranked records; these are discovery results, not 175 verified studies.
Wiley Scholar Gateway was attempted but repeatedly returned an internal error.
Wiley-published originals accessible through other routes are identified normally.
This is a targeted, source-verified review, not an exhaustive systematic review.
Detailed evidence and connector logs accompany the source in
\code{external\_calibration\_review\_evidence/}.

\textbf{Period, population and restrictions.}
Let $\Delta$ denote years per V6 young-to-old transition; this review does not choose it.
Dimensionless risk curvature, markups and static substitution elasticities are not
annualized. Discount probabilities and innovation coefficients need explicit time bridges.
US listed firms, US households, selected foreign countries and the full RoW are distinct
populations. In particular, the current $\rho_{US}>1$ implementation and the sufficient
theorem's $\gamma<1$ must not determine which estimates enter the evidence set.
Solver-excluded candidates such as $\rho_{US}\le1$ require an alternative validated
implementation. The solver accepts positive $\gamma$; values above one instead lose
the cited sufficient theorem and require separate equilibrium checks, without necessarily
requiring a solver change.


\clearpage
\section[Preference weight: M01, beta (Group B)]{Preference weight: M01, $\beta$ (Group B)}

In V6, $\beta$ is the old-age utility weight in
\[
 (1-\beta)\log c_t^y+
 \frac{\beta}{1-\gamma}\log E_t[(c_{t+1}^o)^{1-\gamma}],
 \qquad A_{US,t}=\beta e_{US,t}.
\]
It is a two-period saving weight. It is neither an annual discount factor nor the
present-bias parameter sometimes also called $\beta$.
For two consumption dates with unnormalized weights $1,d$, normalization gives
$\beta=d/(1+d)$. Setting $d=\delta_a^\Delta$ additionally assumes that one V6
transition represents $\Delta$ years and that endpoint utility appropriately
represents both life stages. Aggregating annual utility within life stages can
require a different bridge.

\begin{evidence}
Hirano, Kishi and Toda, pinned 2025 version \src{HKT25}.
Same two-period utility-weight interpretation; Figure~1 note.
&
$\beta=0.5$: theoretical numerical illustration, not an estimate or empirical
range.
&
Analytical overlapping-generations bubble model and numerical figure.
A model period is a generation; the figure does not estimate its calendar length.
&
Strong definition match, absent empirical discipline. Useful for reproducing
the original mechanism; do not label it a literature-estimated baseline.
\\[5pt]
Gourinchas and Parker (2002) \src{GP02}.
Source $\delta$ discounts annual expected CRRA utility; Table~III, p.~71.
&
Estimated $\delta=0.9598$ (corrected SE $0.0179$); optimal-weight variant
$0.9569$ (SE $0.0150$). These are distinct estimators, not interval endpoints.
&
Simulated-moments lifecycle consumption model with earnings risk, liquidity
constraints and retirement. US CEX 1980--1993, with PSID income dynamics;
annual periods.
&
Conditional endpoint bridge:
$\beta=\delta^\Delta/(1+\delta^\Delta)$.
For $\Delta=30$, these points imply $0.2260$ and $0.2105$ (our calculations).
Different EIS, income timing and household budgets make numerical transfer weak.
\\[5pt]
Laibson et al.\ (2026) \src{LLMRT26}.
Annual discount function; Table~3, frequency in footnote~25.
&
Exponential estimate $\delta=0.9601$ (SE $0.0053$).
Present-biased model: continuation $\delta=0.9891$ (SE $0.0051$) and
present-bias weight $b=0.5305$ (SE $0.1143$).
&
Simulated moments from US SCF 1995--2013; high-school households without
college degrees, with credit cards. Annual lifecycle model; 16 moments for
ages 21--60.
&
The exponential point yields conditional $\beta=0.2277$ at $\Delta=30$.
The present-biased continuation factor is not a clean exponential alternative;
$b\neq\beta_{\mathrm{V6}}$. Do not import either parameter directly.
\\[5pt]
\end{evidence}

\noindent\textbf{Recommendation, distinct from evidence.}
Leave the empirical V6 baseline unassigned until life-stage duration and utility
aggregation are specified. Retain $0.5$ solely as the HKT comparison case.
With the explicitly assumed endpoint bridge and $\Delta=20$--$40$, the two
exponential studies' $\delta=0.9569$--$0.9601$ imply approximately
$\beta=0.1465$--$0.3070$. Adding the present-biased continuation factor as a
separate, qualified scenario extends the upper endpoint to $0.4454$.
These are calculated scenarios, not confidence intervals or evidence that V6
must satisfy them. Mapping confidence: high conditional algebra, low empirical
transferability.

\clearpage
\section[Risk curvature: M02, gamma (Group A)]{Risk curvature: M02, $\gamma$ (Group A)}

V6's old-age certainty equivalent is
$\{E_t[(c_{t+1}^o)^{1-\gamma}]\}^{1/(1-\gamma)}$.
Thus $\gamma$ controls risk aversion, while the intertemporal elasticity of
substitution (EIS) is fixed at one. Its numerical value is not annualized.
The sufficient-theorem condition $\gamma<1$ is a model restriction, not a
criterion for accepting evidence. Values above one remain relevant candidates;
the displayed formula needs its limiting interpretation at $\gamma=1$.

\begin{evidence}
Gourinchas and Parker (2002) \src{GP02}.
Source CRRA curvature; Tables~III and V, pp.~71,79.
&
Estimated curvature $0.5140$ (SE $0.1707$) or $1.3969$ (SE $0.1137$), depending
on weighting. Specification variants span $0.1278$--$5.2823$; some imply
implausible initial wealth.
&
Annual lifecycle simulated moments; US CEX 1980--1993 and PSID income risks.
Risk curvature and annual discounting are jointly estimated.
&
Risk-curvature interpretation is comparable, but source EIS equals inverse
CRRA, whereas V6 separates risk aversion and fixes EIS at one. Neither the
specification span nor selected low estimates constitute a V6 confidence range.
\\[5pt]
Laibson et al.\ (2026) \src{LLMRT26}.
Source consumption-utility CRRA curvature; Table~3.
&
Estimated $1.4663$ (SE $0.2256$) with exponential discounting;
$1.9355$ (SE $0.4350$) with present bias.
&
Annual US lifecycle model; same SCF population and moments described above.
Curvature is jointly fitted with discount parameters.
&
Supports a low-to-moderate alternative near two. No Epstein--Zin separation;
liquidity constraints and present bias alter effective intertemporal responses.
Do not identify effective EIS mechanically with inverse curvature in the
present-biased model.
\\[5pt]
Barsky et al.\ (1997) \src{BJKS97}, Table~XI, p.~563;
Kimball et al.\ (2008) \src{KSS08}, Table~4.
CRRA over lifetime resources inferred from hypothetical income gambles.
&
BJKS: mean risk aversion $12.1193$, harmonic mean about $4.2$.
KSS, correcting response error: mean $8.2$, median $6.3$, interquartile range
$3.9$--$10.3$, 5th--95th percentiles $1.9$--$20.8$.
&
US Health and Retirement Study, older respondents and spouses.
KSS uses repeated 1992--2002 questions, response-error and status-quo
adjustments. This is lifetime-risk elicitation, not an annual asset-pricing model.
&
The ranges describe population heterogeneity, not estimation confidence.
A representative V6 curvature requires an aggregation argument.
The reciprocal of mean risk tolerance is a harmonic risk-aversion statistic,
not the mean; do not silently use it as representative $\gamma$.
\\[5pt]
Holt and Laury (2002) \src{HL02}.
CRRA utility of experimental cash payoffs; Table~3 and p.~1649.
&
Low-payoff behavior centers roughly around $0.3$--$0.5$; one
choice-implied interval is $0.41$--$0.68$. Risk aversion increases when
real stakes increase. These intervals are not confidence intervals.
&
Incentivized paired-lottery choices by students and faculty at three US
universities; one-shot experiments rather than a calendar-frequency model.
&
Evidence supports retaining subunit-curvature scenarios. Transferability is
weak: small cash-payoff risk, stakes dependence and a selected subject pool
differ from V6 old-age consumption risk.
\\[5pt]
\end{evidence}

\clearpage
\subsection{Asset-pricing estimates, synthesis and proposed risk-curvature scenarios}
\begin{evidence}
Vissing-Jorgensen and Attanasio (2003) \src{VA03}.
Epstein--Zin risk curvature; author-uploaded January~10 extended manuscript,
Table~2A, case~4, p.~14.
&
All-stockholder estimates $8.1,6.3,5.5$ as assumed human-wealth share rises
from $1/3$ to $2/3$ to $0.9$; richest-third estimates $28.4,19.9,16.5$.
Middle all-stockholder case: 95\% interval $[0.3,13.9]$.
&
US CEX 1982--1996 stockholders; semiannual consumption growth with overlapping
monthly observations, asset returns and VAR-based human-wealth adjustments.
Structural asset-pricing estimation.
&
Conceptually closer because risk curvature and EIS are separated. Numerical
mapping remains weak: EIS, wealth aggregation and human-capital assumptions
differ. Numbers are verified in the dated extended manuscript, not asserted
to be an identically numbered table in the shorter published article.
\\[5pt]
Chen, Favilukis and Ludvigson (2013) \src{CFL2013}.
Source $\theta$ is risk aversion; $\rho^{-1}$ is EIS.
Equations~(1)--(2), p.~44; Table~2, p.~62.
&
Stockholder $\theta=20,17$ with 95\% intervals $[0.25,40]$, $[1,43.3]$;
aggregate $\theta=57.5,60$ with intervals $[27.5,129]$, $[42,144]$.
EIS estimates exceed one.
&
Quarterly US data 1952Q1--2005Q1; consumption, Treasury bills and equity
portfolios. Profile sieve minimum distance with second-step GMM.
Stockholder consumption is a fitted mimicking factor.
&
Map $\gamma=\theta$ conceptually, without annualization. Infinite-horizon
continuation-value risk and freely estimated EIS differ from V6.
Wide bootstrap intervals and possible model misspecification limit transfer.
The span $17$--$60$ is between specifications, not a confidence interval.
\\[5pt]
Elminejad, Havranek and Irsova (2025) \src{EHI25}.
Relative-risk-aversion meta-analysis; Table~6 and main discussion.
&
Bias-adjusted summaries differ by context: roughly one in economics,
two--seven in finance. US-context prediction $5.20$, 95\% interval
$[2.28,8.11]$.
&
Synthesis of 1,021 estimates from 92 studies, with literature search ending
May~2022. Studies differ in populations, frequencies and empirical designs.
This is secondary synthesis, not a new original estimate.
&
Useful context for choosing a sensitivity grid; not a direct V6 measurement.
Its discipline split cautions against pooling lifecycle and asset-pricing
evidence into a single consensus number.
\\[5pt]
\end{evidence}

\noindent\textbf{Recommendation, distinct from evidence.}
Use $\gamma=5$ as a provisional finance-oriented candidate, with $2$ as an
explicit lifecycle-oriented alternative. A core sensitivity grid
$\{0.5,2,5,6.3,10\}$ preserves low experimental and higher survey/finance
evidence; add a separate high-risk family $\{17,20,60\}$ from recursive
asset-pricing estimates. This grid is a research recommendation, not a pooled
credible interval. Conceptual mapping confidence is medium; numerical
confidence is low. The candidate $5$ violates the sufficient-theorem
restriction $\gamma<1$; retain and report that conflict instead of choosing
a low number solely to preserve the theorem.


% Production-parameter chapter. Sources use P1--P12 from production_review.json.
% Parent report supplies the evidence environment and \src macro.
\clearpage
\section[Inverse markups: competing empirical cost definitions]{Inverse markups: competing empirical cost definitions}
\textbf{V6 definition (M10).} For $i\in\{US,W\}$,
\[
X_i=N_i^{1-1/\vartheta_i}\left(\int_0^{N_i}x_i(j)^{\vartheta_i}\,dj\right)^{1/\vartheta_i},
\qquad p_i=\frac{w_{H,i}}{\vartheta_i},\qquad
\mu_i=\frac{p_i}{MC_i}=\frac1{\vartheta_i},\qquad
\epsilon_i=\frac1{1-\vartheta_i}=\frac{\mu_i}{\mu_i-1}.
\]
Each intermediate uses skilled labor one-for-one. Thus $\vartheta_i$ is inverse gross markup; $\epsilon_i$ is the variety demand elasticity, distinct from the final-sector elasticity $1/\rho_i$. These dimensionless quantities require no annual-to-cohort exponentiation.

\begin{evidence}
\src{DLEU20}\par De Loecker, Eeckhout and Unger (2020). Gross price/marginal-cost markup; sales-weighted across firms.
& \textbf{Estimated:} $\mu=1.21$ (1980), $1.61$ (2016). Hence $\vartheta=0.826,\ 0.621$. Dated point estimates, not confidence bounds. Author manuscript, Section 3.1/Figure 1, pp.\ 9--10.
& Annual U.S. Compustat, 1955--2016; sector/time-specific production-function elasticities divided by flexible-input revenue shares. Census comparisons address listed-firm selection. Within-year input flexibility is maintained.
& Exact reciprocal conditional on the cost definition. Economic transfer is qualified: heterogeneous firms, overhead, selection and reallocation differ from V6's homogeneous intermediate sector. $1/\overline\mu\ne\overline{1/\mu}$. Do not also impose empirical net-profit shares without reconciling omitted costs.\\[4pt]

\src{Traina18}\par Traina (2018), working paper. Markup treating operating expenses, including SG\&A, as variable inputs.
& \textbf{Estimated:} approximately $\mu=1.15$ (1950), $1.10$ (1980), $1.15$ (sample end); $\vartheta=0.870,\ 0.909,\ 0.870$. Section ``Longer Time Horizons,'' Figure 5 discussion; original text verified.
& U.S. Compustat, 1950--2016; industry-specific Cobb--Douglas production estimation. Utilities and financials excluded. Annual and three-year production windows produce similar paths.
& Preserve this lower-markup alternative. The disputed variable-versus-fixed treatment of marketing/management costs explains much of the difference from the COGS approach. This is a separate measurement specification, not sampling uncertainty around the preceding estimates. V6's intermediate cost structure still differs.\\[4pt]

\src{DE21}\par De Loecker and Eeckhout, \emph{Global Market Power}, February 2021 version. Regional sales-weighted markups.
& \textbf{Estimated, 2016:} Europe $1.63$, Asia $1.45$, Africa $1.38$, Oceania $1.55$, South America $1.59$; inverse values $0.613,0.690,0.725,0.645,0.629$. U.S. $1.84\Rightarrow0.543$. Table 1, p.\ 13.
& Annual 1980--2016; 67,491 firms in 134 countries. Global listed-firm accounts, with U.S.-based industry output elasticities. Sales-weighted aggregation; headquarter location need not equal production location.
& Useful geographic alternatives, not a full-RoW coefficient. The U.S. estimate differs from the preceding paper under a different database/specification. The global aggregate includes the U.S. Define the RoW universe and aggregation weights before choosing $\vartheta_W$.\\[4pt]
\end{evidence}

\textbf{Methodological disagreement.} \src{BHKZ21} Bond et al.\ (2021) question substituting revenue elasticities for output elasticities and treating imperfectly flexible or demand-shifting inputs as flexible. \src{DL21} De Loecker (2021) emphasizes additional price/quantity data or structural assumptions as remedies. Neither gives a new numerical V6 coefficient; the debate lowers confidence in a mechanical aggregate-markup transfer.

\textbf{Recommendation, separate from evidence.} For a modern U.S. listed-firm mapping, retain $\vartheta_{US}=0.621$ and $0.870$ as competing cost-definition candidates. Use $[0.621,0.909]$, extended to $0.543$, as a measurement-specification envelope, not a confidence interval. Defer a single RoW point; Europe $0.613$ and Asia $0.690$ are explicit geographic anchors. Algebraic mapping confidence is high; economic transfer is low to medium for the U.S. and low for RoW.

\clearpage
\section[Final-sector curvature: the skill-substitution evidence]{Final-sector curvature: the skill-substitution evidence}
\textbf{V6 definition (M12).} With $\widetilde X_i=A_{X,i}X_i$ and $\widetilde L_i=A_{L,i}L_i$,
\[
F_i=\left[\alpha_i\widetilde X_i^{1-\rho_i}+(1-\alpha_i)\widetilde L_i^{1-\rho_i}\right]^{1/(1-\rho_i)},
\qquad \sigma_i=\frac1{\rho_i},\qquad X_i=\varphi_i H_i.
\]
The relevant margin is skilled labor allocated to intermediate production versus unskilled labor. Total skilled labor includes R\&D. The current U.S. solver requires $\rho_{US}>1$, hence $\sigma_{US}<1$; retain evidence outside this restriction.

\begin{evidence}
\src{KM92}\par Katz and Murphy (1992). College/high-school-equivalent labor substitution.
& \textbf{Estimated:} inverse-demand slope $-0.709$ (SE $0.150$); $\sigma\simeq1.41$. Equation (19), p.\ 69; discussion p.\ 72. Their plotted $\sigma=0.5,4$ alternatives are scenarios, not estimated bounds.
& Annual U.S. CPS, 1963--1987; efficiency-unit labor aggregates. OLS relative wages on relative supplies and a linear technology trend, interpreted through competitive two-labor CES demand.
& Conditional proxy $\rho=0.709$. Skill definitions, technology trends, V6 markups and the production/R\&D split require alignment. Trend-controlled OLS is not independent identification of the technology process. This candidate cannot run under the current $\rho_{US}>1$ restriction.\\[4pt]

\src{CP05}\par Ciccone and Peri (2005). Long-run substitution between education groups.
& \textbf{Estimated:} preferred $\sigma=1.50$ (SE $0.44$); IV variants $1.20$--$2.00$; fixed-effects OLS $2.85$--$3.44$. Table 4, p.\ 656. These estimator ranges are not confidence intervals.
& Five decennial censuses, 1950--1990, 48 U.S. states. Schooling/child-labor-law instruments, state/year effects; preferred Fuller LIML. Wages: U.S.-born white full-time males aged 40--50; baseline supply: white males 21--59.
& Preferred $\rho=0.667$; IV point-sensitivity $0.500$--$0.833$. Long-run education substitution is a qualified skill-input proxy, not an exact V6 or RoW estimate. Changing schooling aggregation and weak-instrument treatment matters.\\[4pt]

\src{KORV00}\par Krusell, Ohanian, R\'{i}os-Rull and Violante (2000). Distinct nested-CES elasticities.
& \textbf{Estimated:} outer unskilled-versus-equipment/skilled elasticity $1.67$; inner equipment-versus-skilled elasticity $0.67$. Table II, p.\ 1041. Alternative skill threshold gives outer elasticity $1.89$ (p.\ 1044).
& Annual U.S. CPS and capital/quality-adjusted equipment-price data, 1963--1992. Two-step simulated pseudo-maximum likelihood; structures, equipment, two labor types and latent labor efficiencies.
& Outer-nest analogy: $\rho=1/1.67=0.599$. Inverting the inner $0.67$ to obtain $1.493$ maps the \emph{wrong input pair}; it cannot validate V6 complementarity. The complete nested technology differs from V6.\\[4pt]

\src{HILZ24}\par Havranek, Irsova, Laslopova and Zeynalova (2024 issue; online 2022). Negative inverse skill elasticity.
& \textbf{Meta-estimated:} Table 5, p.\ 1197: all countries $-0.27$ (95\% credible interval $[-0.48,-0.05]$); U.S. $-0.16$ $[-0.38,0.06]$; developing $-0.47$ $[-0.70,-0.24]$. Implied $\sigma=3.7,6.3,2.1$.
& 682 estimates/77 studies; publication- and attenuation-bias analysis, model averaging. Mixed populations/frequencies; preferred profile favors annual panels, multilevel CES, fixed effects and valid IV.
& Skill-proxy $\rho=0.27,0.16,0.47$. U.S. mapped interval $[-0.06,0.38]$ crosses zero; do not report a finite positive elasticity interval. Bias-correction/IV assumptions matter. This newer evidence challenges a universal $\sigma=1.5$ benchmark.\\[4pt]
\end{evidence}

\textbf{Convention and version audit.} \src{CC06} Caselli--Coleman (2006, pp.\ 504--505) \emph{borrow} $\sigma=1.4$ and use $1.1,1.4,1.7,2$ for robustness; these are not independent estimates. Their CES exponent $s$ maps as $\rho_{V6}=1-s=1/\sigma$. The 2020/2021 ``Less Substitutable'' working-paper version of \src{HILZ24}, reporting $0.6$--$0.9$, is superseded by the final 2024 paper and must not count as independent evidence.

\clearpage
\section[Complementarity evidence and the unresolved V6 input mapping]{Complementarity evidence and the unresolved V6 input mapping}
The following papers support complementarity of \emph{physical capital and labor}. They deserve consideration alongside skill-substitution evidence, but their input pairs differ. A theorem-compatible reciprocal is not sufficient evidence for the V6 coefficient.

\begin{evidence}
\src{OR21}\par Oberfield and Raval (2021). Aggregate physical capital--labor elasticity, including reallocation.
& \textbf{Estimated:} aggregate manufacturing $\sigma=0.5$--$0.7$; plant-level $0.3$--$0.5$ (pp.\ 703, 705). Figure 4b, p.\ 719: $0.60\to0.54$ over 1972--2007 with micro elasticities fixed at 1997 values.
& U.S. Census of Manufactures/Annual Survey of Manufactures. Local-wage IV identifies plant elasticities; amenities and shift-share instruments; heterogeneity determines aggregation. Annual/census observations, not a cohort parameter.
& Reciprocal \emph{analogy}: $\rho=1.43$--$2.00$, with $1.67$ from $\sigma=0.60$. The source does not impose a stable aggregate CES. Physical $K$ versus total labor is not V6's $X=\varphi H$ versus $L$; transfer confidence is low.\\[4pt]

\src{GHIK22}\par Gechert, Havranek, Irsova and Kolcunova (2022). Conditional meta-estimates of capital--labor substitution.
& \textbf{Meta-estimated:} Table 9, p.\ 77: best practice $\sigma=0.30$, 95\% CI $[-0.01,0.60]$; country-data variant $0.50$ $[0.18,0.81]$; quarterly variant $0.42$ $[0.08,0.76]$.
& 3,186 estimates/121 studies published before 2019. Nonlinear publication-bias corrections, Bayesian/frequentist model averaging, 71 study characteristics. Mixed countries, sectors and frequencies; methodological profile is consequential.
& Point reciprocals $3.33,2.00,2.38$ are capital--labor analogies. The baseline interval crosses zero: it cannot become a bounded positive $\rho$ interval by inversion. Neither the synthetic-study mean nor its methodological choices directly identify the V6 input margin.\\[4pt]
\end{evidence}

\subsection{What must be held fixed in the empirical mapping}
V6's factor prices and symmetry imply
\[
\log\frac{w_{H,i}}{w_{L,i}}
=\log\frac{\vartheta_i\alpha_i}{1-\alpha_i}
+(1-\rho_i)\log\frac{A_{X,i}}{A_{L,i}}
-\rho_i\log\frac{X_i}{L_i},\qquad X_i=\varphi_iH_i.
\]
Consequently, a regression on total $H_i/L_i$ need not recover $-\rho_i$ when R\&D allocation, knowledge-driven augmentations or markups move with relative labor supply. The two literatures also disagree about the relevant input pair and aggregation. Their results should remain separate empirical specifications; averaging them into one ``consensus'' curvature would discard the economic distinction.

\subsection{Proposed candidates and sensitivity design}
\textbf{Skill-input branch.} Retain $\rho=0.667$ (Ciccone--Peri) and $0.709$ (Katz--Murphy) as original-study candidates, alongside the newer meta-evidence candidates $0.27$ (all countries), $0.16$ (U.S.) and $0.47$ (developing countries). An initial \emph{analyst} scenario envelope is $[0.16,0.91]$, with $0.05$ as a high-substitution stress point suggested by the all-country meta interval. These are not a pooled estimate or V6 confidence interval. RoW application requires explicit geographic/skill pooling assumptions.

\textbf{Capital--labor analogy branch.} Consider $\rho=2$ as a separate complementarity candidate, and $[1.43,3.33]$ as a broad scenario envelope drawn from the two studies above. The familiar HKT illustration at $\rho=2$ does not turn this analogy into an estimated V6 benchmark. Keep low mapping confidence visible.

\textbf{Implementation implication.} The central skill-proxy candidates cannot run under the current $\rho_{US}>1$ restriction. Report this as a conflict between an empirical mapping and the implemented model class; do not discard the estimates or silently truncate them. Any operational choice of the complementarity branch remains conditional on defending its economic mapping. Elasticities stay dimensionless when changing the model period.


\clearpage
\section[Regime persistence: M07, pi (Group B)]{Regime persistence: M07, $\pi$ (Group B)}

V6 defines $\pi=\Pr(z_{t+1}=u\mid z_t=u)$, with $b$ absorbing.
An empirical transition therefore needs both an operational definition of the
unbalanced-to-balanced technology switch and a model-period length $\Delta$ in
years. A recurring productivity regime is the closest analogy reviewed below;
neither a recession nor an equity-price fall defines the V6 event.

\begin{evidence}
Kahn and Rich (2007) \src{KR07}.
Quarterly high-growth persistence $q_1=\Pr(S_{1t}=0\mid S_{1,t-1}=0)$;
equations~(7)--(9), p.~1677; Table~2, p.~1683.
&
Estimated $q_1=0.990$ (SE $0.011$), implying a 100-quarter mean spell.
Low-growth persistence is $0.983$ (SE $0.015$).
Alternative specifications give high-growth persistence $0.991$ or $0.993$;
these are not confidence bounds.
&
Approximate maximum-likelihood Markov dynamic-factor model; US quarterly
1947Q1--2002Q4, December~2005 vintage. Output/hour, compensation/hour,
consumption/hour and detrended hours identify common trend growth. Both
growth regimes recur.
&
Conditional first-exit analogy: $\pi_\Delta=q_1^{4\Delta}$.
This is uninterrupted survival, not the recurrent chain's endpoint high-state
probability. Aggregate growth regimes differ from V6 factor-spillover regimes.
Per-capita specification loses regime identification (p.~1687);
few transitions limit precision.
\\[5pt]
Barro (2006) \src{Barro06}.
Poisson intensity $p$ of downward output jumps; pp.~825,831 and Table~V, p.~846.
&
Data-based frequency calibration: $p=0.017$ per year, with reported
frequency range $0.015$--$0.020$.
The $0.025$ sensitivity case is calibrated, not an estimated upper bound.
&
60 cumulative per-capita GDP contractions of at least 15\%, 35 countries,
approximately the twentieth century. Recurring rare disasters in a
representative-agent asset-pricing economy.
&
The source's no-jump probability is $\exp(-p\Delta)$.
Equating it to V6 $\pi$ additionally equates two different events.
Useful disaster-risk sensitivity only; a global GDP-contraction frequency
does not estimate an absorbing US technology transition.
\\[5pt]
Barro and Urs\'ua (2008) \src{BU08}.
Annual normal-to-disaster transition probability $h$; Tables~10--11,
pp.~296--297.
&
At a 10\% contraction threshold, $h=0.0363$ for consumption and $0.0369$
for GDP. At 15\%, corresponding values are $0.0218$ and $0.0192$.
These event-definition alternatives are not a confidence interval.
&
Historical annual series starting as early as 1870; 24 consumption and
36 GDP countries. Events are peak-to-trough cumulative contractions;
hazards use normal-state exposure years.
&
Under a matched constant annual hazard, $\pi=(1-h)^\Delta$.
Events recur and recover. The source's symbol $\pi$ instead denotes
disaster exit and must not be copied into V6. Published sample counts
differ from the older working-paper abstract.
\\[5pt]
\end{evidence}

\noindent\textbf{Recommendation, distinct from evidence.}
No directly estimated V6 benchmark is available. Prefer the Kahn--Rich
\emph{conditional duration scenario}, $\pi=0.990^{4\Delta}$, to a generic
disaster center: it implies $0.9606$ for one year and $0.3660$ for 25 years
(our calculations). A 10/25/50-year mean-duration stress design can use
$q_D=1-1/(4D)$ and $\pi=q_D^{4\Delta}$; only the 25-year center is
source-derived, while the endpoints are analyst choices. This is not a
credible interval: the reported Wald interval for $q_1$ crosses one, making
inverse-duration inference unreliable. Mapping confidence is low for the
V6 event, despite exact conditional frequency conversion. Use the published
2007 estimates; 2003 drafts differ in notation and sample endpoints.

\clearpage
\section[Innovation productivity: M09, a US ,a W (Group B)]{Innovation productivity: M09, $a_{US},a_W$ (Group B)}

V6 imposes $N_{i,t+1}/N_{i,t}-1=a_i(1-\varphi_{i,t})H_i$.
For matched gross knowledge growth $G_{N,i}$ and research share
$s_{R,i}=1-\varphi_i$, conditional recovery and frequency conversion are
\[
 a_iH_i=\frac{G_{N,i}-1}{s_{R,i}}
       =\frac{(1+g_{N,i})^\Delta-1}{s_{R,i}}.
\]
The second equality assumes constant annual knowledge growth $g_{N,i}$ and
research share over $\Delta$ years; it matches endpoint knowledge growth,
not the full temporally aggregated OLG equilibrium. Labor units matter:
if $H_i'=cH_i$, then $a_i'=a_i/c$.
Recovery requires $s_{R,i}>0$. A zero-research observation has $G_{N,i}=1$
and cannot identify $a_i$; share uncertainty approaching zero can make
an upper bound unbounded.

\begin{evidence}
Hirano, Kishi and Toda, arXiv v2 (2025) \src{HKT25}.
Same expanding-variety research law; Figure~1, p.~24.
&
$aH=0.2$: numerical illustration without an estimated population or
empirical range.
&
Closed-economy two-period OLG bubble model. The figure does not establish
the calendar length of a model period or estimate a country-specific
research coefficient.
&
Strong definition match for the product $aH$; no empirical level transfer.
The value is not 0.2 annually and does not establish equal US and RoW
productivity. Recover $a_i$ only after fixing labor units.
\\[5pt]
Bloom et al.\ (2020) \src{BJRW20}.
Research productivity equals idea-output growth divided by effective
researchers; equations~(1)--(3), pp.~1108--1110; Table~7, p.~1134.
&
Reported annual productivity growth: aggregate US $-5.1\%$;
Moore's Law $-6.8\%$; cotton, version~2, $+1.3\%$.
These are descriptive growth-accounting ratios, not estimated levels of $a_i$.
&
US aggregate 1930--2015: decadal TFP growth and intellectual-property
investment deflated by skilled wages. Moore's Law: 1971--2014;
cotton: 1969--2009. Input quality and output concepts vary by application.
&
Declining research productivity challenges constant V6 $a_i$, while cotton
preserves a contrary case. TFP or chip performance does not directly count
V6 varieties. Matched knowledge, labor units, quality and period are needed;
never insert these growth rates as $a_i$.
\\[5pt]
Bottazzi and Peri (2007) \src{BP07}.
Patent flow $Pat=\iota RD^\lambda A^\phi A_{ROW}^{\xi}$;
equations~(5)--(9), pp.~492--494; Table~4, p.~501.
&
Estimated long-run elasticities
$(\lambda/(1-\phi),\xi/(1-\phi))=(0.786,0.557)$;
with country trends, $(0.304,0.168)$.
Knowledge depreciation of 10\% is imposed, not estimated.
&
Dynamic-OLS panel cointegration; 15 OECD countries, annual research
personnel and citation-weighted US patent applications.
Text dates are 1973--1999; Table~4 says 1972--1999.
Country effects and trends materially change estimates.
&
V6 instead imposes unit research and own-knowledge exponents, no foreign
knowledge term and no depreciation. The finite cointegrating ratios do
not identify its coefficient $a_i$; source $\xi$ is also not V6 $\xi_u$.
Useful specification discipline, not a portable prior.
\\[5pt]
Jones (1995) \src{Jones95}.
R\&D-based growth and scale effects; publisher abstract verified.
&
Qualitative evidence against standard scale effects; no numerical
coefficient is borrowed from this source.
&
Industrial-economy time-series critique and semi-endogenous growth
theory. Full-text coefficient and equation claims are not relied upon.
&
Motivates a diminishing-productivity specification check. It does not
identify a V6 variety-growth coefficient or resolve its research-labor
normalization.
\\[5pt]
\end{evidence}

\noindent\textbf{Recommendation, distinct from evidence.}
Recover each country's $a_iH_i$ from matched knowledge growth and research
share; leave empirical scalar baselines unset until that bridge exists.
Use $0.2$ only as an HKT initialization, not a portable range.
Propagate jointly credible growth/share bounds through the displayed
formula, separately for US and the declared RoW aggregate; include
knowledge-proxy and declining-productivity specifications. Empirical mapping
confidence is low. \textbf{Source caution:} BJRW Table~7 reports $-5.1\%$
and a 14-year half-life, while its preceding prose reports $-5.3\%$ and
13 years; the table values are retained without resolving that discrepancy.

\clearpage
\section[Knowledge augmentation elasticities: M16, xi u,nu u (Group B)]{Knowledge augmentation elasticities: M16, $\xi_u,\nu_u$ (Group B)}

The exact V6 definitions are
\[
 A_{X,US}=\bar A_{X,US,u}N_{US}^{\xi_u},\qquad
 A_{L,US}=\bar A_{L,US,u}N_{US}^{\nu_u},\qquad
 (\xi_u,\nu_u)=
 \left(\frac{\partial\log A_X}{\partial\log N},
       \frac{\partial\log A_L}{\partial\log N}\right).
\]
These are dimensionless factor-augmentation elasticities, not annual
technology-growth rates, substitution elasticities, or exponents in an
idea-production function. The same knowledge concept must enter both ratios.

\begin{evidence}
Hirano, Kishi and Toda, arXiv v2 (2025) \src{HKT25}.
Factor-augmentation powers; Figure~1, p.~24.
&
$(\xi_u,\lambda_u)=(0.7,0.2)$ is an illustrative calibration;
source $\lambda_u$ maps to V6 $\nu_u$. No estimated joint range is supplied.
&
Theoretical closed-economy OLG numerical example, with unspecified calendar
date length. Its positive post-switch exponents differ from the imposed
zero-exponent specialization in V6.
&
Same pre-switch definition permits mechanism comparison without
annualization. It provides no empirical basis for a tight neighborhood
around $(0.7,0.2)$ or for selecting only values admitted by the V6 solver.
\\[5pt]
Katz and Murphy (1992) \src{KM92}.
Log college/high-school wage ratio regressed on relative supply and time;
equation~(19), p.~69.
&
Estimated supply slope $-0.709$ (SE $0.150$); annual trend $0.033$
(SE $0.007$). The trend is a conditional wage-demand shift, not
$\xi_u-\nu_u$ or either individual augmentation elasticity.
&
OLS on annual US CPS wage and labor-supply data, 1963--1987;
competitive two-skill CES interpretation. No separately measured
V6 variety stock or research allocation enters this regression.
&
V6 uses an intermediate bundle, a markup and endogenous skilled-labor
allocation between production and research. Its wage mapping below
requires all three plus knowledge growth. In particular, the sign of
the augmentation term changes with $1-\rho$.
\\[5pt]
Bloom et al.\ (2020) \src{BJRW20}, equation~(17) and Table~7, p.~1134;
Bottazzi and Peri (2007) \src{BP07}, equations~(5)--(9).
&
BJRW's research-productivity equation is
$\dot A/A=(\alpha A^{-\beta})S$:
reported $\beta=3.1$ for aggregate output and $0.2$ for Moore's Law.
BP estimates the distinct long-run ratios reported above.
&
BJRW uses aggregate and application-specific growth accounting;
BP uses an annual international patent/R\&D panel. Neither estimates
two final-production augmentations against a common variety stock.
&
BJRW $\beta$ governs diminishing research productivity, not V6 saving
or augmentation. BP $\xi$ governs foreign knowledge in patent production.
These numerical coefficients cannot supply values or ranges for
V6 $\xi_u,\nu_u$ despite related terminology.
\\[5pt]
\end{evidence}

\noindent\textbf{Exact conditional mapping derived from V6.}
With fixed markup parameter $\vartheta$ and CES weight $\alpha$,
\[
 \frac{w_H}{w_L}=\vartheta\frac{\alpha}{1-\alpha}
 \left(\frac{A_X}{A_L}\right)^{1-\rho}
 \left(\frac{\varphi H}{L}\right)^{-\rho},\qquad
 \Delta\log(w_H/w_L)+\rho\Delta\log(\varphi H/L)
 =(1-\rho)(\xi_u-\nu_u)\Delta\log N.
\]
Thus relative wages can discipline a combination only conditionally;
they do not identify both exponents. Separate augmentation-growth evidence
is needed for $\xi_u=\Delta\log A_X/\Delta\log N$ and
$\nu_u=\Delta\log A_L/\Delta\log N$.

\noindent\textbf{Recommendation, distinct from evidence.}
Retain $(0.7,0.2)$ solely as the HKT comparison case; leave the empirical
baseline and joint range unassigned. Construct a joint growth-ratio region
from matched augmentation and knowledge evidence, retaining zero, equal or
reversed augmentation growth if supported, even when the current solver
excludes it. Multiplying $N$ by a constant changes scales, not exponents;
redefining knowledge by a power changes exponents. Confidence is high in
the HKT definition match and low in empirical numerical transfer.


\clearpage
\section[Distribution weights and technology levels: M11, M13--M15 (Group B)]{Distribution weights and technology levels: M11, M13--M15 (Group B)}

The authoritative V6 definitions are
\[
Y_i=\left[\alpha_i(A_{X,i}X_i)^{1-\rho_i}
 +(1-\alpha_i)(A_{L,i}L_i)^{1-\rho_i}\right]^{1/(1-\rho_i)},
\qquad X_i=\varphi_iH_i .
\]
M11 comprises $\alpha_{US},\alpha_W$. M13 comprises
$A^u_{X,US}=\bar A_{X,US,u}N_{US}^{\xi_u}$ and
$A^u_{L,US}=\bar A_{L,US,u}N_{US}^{\nu_u}$.
M14 comprises $A^b_{X,US}=\bar A_{X,US,b}$ and
$A^b_{L,US}=\bar A_{L,US,b}$ because $\nu_b=0$.
M15 comprises $A_{X,W}=\bar A_{X,W}$ and $A_{L,W}=\bar A_{L,W}$ because $\xi_W=0$.
These six barred coefficients are augmentation scales, not observed aggregate TFP.

\begin{evidence}
Caselli and Coleman (2006) \src{CC06}.
Skill-specific augmentations $A_s,A_u$ in a labor CES inside a capital--labor technology;
equations (2)--(3), pp.~499--502.
&
Baseline skill elasticity $\sigma=1.4$ and capital share $1/3$ are imposed.
Augmentations are jointly recovered from output and the skill premium.
There is no unit-free common level or range for V6's barred coefficients.
&
Cross-country development accounting: 52 countries; output/capital 1988 and
schooling 1985. Baseline skill split is completed primary schooling.
Static country comparison; PPP quantities and wage-premium inputs.
&
Transfer the inversion method, not the source capital share to V6 $\alpha$.
Adjust for V6's intermediate markup, production-only skilled labor and omitted physical
capital. Cross-country technology differences are not time-series elasticities in $N$.
\\[5pt]
Rossi (2022) \src{Rossi22}.
Relative skill efficiency $(A_HQ_H)/(A_LQ_L)$; equations (4)--(8), Table~I
in the July~2021 accepted manuscript, pp.~7--8,31.
&
At imposed $\sigma=1.5$, US relative efficiency is normalized to 1;
Canada is $0.711$, India $0.041$. These are country inversions,
not a confidence interval or V6 RoW scale range.
&
Comparable wage and hours microdata for 12 countries around 2000.
Some tertiary education defines high skill.
Separate migrant evidence distinguishes embodied human capital $Q$ from technology.
Static CES, competitive labor pricing.
&
Useful for relative-productivity measurement. The reported object mixes
technology and worker quality, and lacks V6's research-labor and markup adjustments.
Re-invert the chosen aggregate; averaging these country numbers does not produce
$(\bar A_{X,W},\bar A_{L,W})$.
\\[5pt]
Cantore and Levine (2012) \src{CL12};
Klump, McAdam and Willman (2012) \src{KMW12}.
CES distribution and scale normalization.
&
No independent empirical point estimate for V6 $\alpha$ or augmentation levels.
Distribution coefficients depend on units and the normalization point.
Normalized baseline income shares are different objects from an unrestricted raw CES weight.
&
Methodological macroeconomic analyses; RBC examples and a theory/empirics survey.
Detailed checks use the 2011 Dynare and ECB working-paper versions,
respectively. No common population or sample-based confidence interval.
&
For $\rho\ne1$, fix one convention and jointly recover the remaining scale
coefficients. Varying $\rho$ while holding all raw scales fixed changes
baseline shares as well as substitution. Re-normalize around declared production moments.
\\[5pt]
León-Ledesma, McAdam and Willman (2010) \src{LMW10}.
Joint production/first-order-condition system with biased technical change.
&
Identification-method evidence, not a new portable value for V6 technology.
Normalized system estimation can improve separation of substitution and factor bias
relative to a single production equation.
&
Monte Carlo comparisons of CES estimators under specified data-generating processes.
Published AER article; full-text method checked in ECB working paper~1001 (2009).
&
Supports joint treatment of $\rho$, shares and technology paths.
Does not establish identification in V6, whose monopoly wedge and endogenous research
allocation introduce additional conditions.
\\[5pt]
Hirano, Kishi and Toda (2025) \src{HKT25}, Figure~1 note.
Same CES weight and factor-augmentation interpretation.
&
Illustration: $\alpha=0.5$, $\rho=2$, $A_XH=10$, $A_LL=1$, $n_0=1$.
Thus $A_X=10/H$ and $A_L=1/L$ in those units, not universal scale values $10$ and $1$.
&
Closed-economy theoretical figure with a generational period;
no fitted population or sampling uncertainty.
Its post-switch exponents are positive, unlike the V6 zero-exponent case.
&
Useful numerical comparison only. The figure does not separately discipline
US/RoW or the two constant V6 post-switch scales.
Do not reinterpret this illustration as evidence for a narrow productivity range.
\\[5pt]
\end{evidence}

\clearpage
\subsection{Explicit V6 mapping and proposed assignment}

\textbf{Our derivation from the V6 equations.}
For each country and reference observation, suppress country subscripts and set
$X_0=\varphi_0H$. Define the intermediate expenditure share, including its monopoly profits,
\[
s_{X,0}\equiv\frac{P_0X_0}{Y_0}
=\frac{w_{H,0}\varphi_0H}{\vartheta Y_0},
\qquad 1-s_{X,0}=\frac{w_{L,0}L}{Y_0}.
\]
The first share is \emph{not} total skilled payroll over output: research payroll is excluded,
while profits on current intermediate production are included through $1/\vartheta$.
Matched observations must satisfy these two shares adding to one under the maintained
production accounting. If they fail, report the scope/measurement residual rather than
treating inconsistent share estimates as independent exact restrictions.
Given $0<s_{X,0}<1$, fixed $\alpha\in(0,1)$ and $\rho\ne1$,
\begin{align}
B_X(\alpha,\rho;s_{X,0},Y_0,X_0)
&\equiv \frac{Y_0}{X_0}
 \left(\frac{s_{X,0}}{\alpha}\right)^{1/(1-\rho)}
= A_{X,0}, \label{eq:bx}\\
B_L(\alpha,\rho;s_{X,0},Y_0,L)
&\equiv \frac{Y_0}{L}
 \left(\frac{1-s_{X,0}}{1-\alpha}\right)^{1/(1-\rho)}
= A_{L,0}. \label{eq:bl}
\end{align}
These are conditional production inversions, not an independent demonstration that
research free entry and the complete equilibrium support the chosen $\varphi_0$.

\begin{evidence}
M11: $\alpha_{US},\alpha_W$.
&
\textbf{Proposed convention:} $\alpha_i=0.5$ with joint rescaling of augmentations,
when $\rho_i\ne1$. No empirical confidence interval is assigned.
&
Only $\alpha_iA_{X,i}^{1-\rho_i}$ and
$(1-\alpha_i)A_{L,i}^{1-\rho_i}$ enter the raw CES.
An alternative $\alpha_i'$ is equivalent only if both combinations are preserved.
&
This is a normalization, not a half-income-share finding.
At $\rho=1$, $\alpha$ becomes the Cobb--Douglas share and cannot be varied
as a pure scale convention. Near one, the inversion is ill-conditioned.
\\[5pt]
M13: $\bar A_{X,US,u}$ and $\bar A_{L,US,u}$.
&
\textbf{Proposed baseline:}
$B_{X,US}/N_{US,0}^{\xi_u}$ and
$B_{L,US}/N_{US,0}^{\nu_u}$.
No portable numerical level is assigned before matching data.
&
Initial US output, production-skilled and unskilled inputs, matched wages,
markup, and a declared $N_{US,0}$ normalization. These production moments
become calibration targets if used here.
&
Propagate wage/share, research-allocation and $\rho$ uncertainty jointly.
Changing units changes the bars. A range imported from a different normalization
would create artificial precision.
\\[5pt]
M15: $\bar A_{X,W}$ and $\bar A_{L,W}$.
&
\textbf{Proposed baseline:} $B_{X,W}$ and $B_{L,W}$ for the declared RoW aggregate.
No numerical full-RoW candidate is presently supported.
&
Matched country coverage, output units, skill definition and wage accounting.
CC06/Rossi provide measurement designs, not a ready-made V6 aggregate.
&
Use alternative coverage and input definitions to construct joint ranges.
Maintain the US--RoW relative scale under any common output-unit normalization.
Covered-country estimates must retain their coverage label.
\\[5pt]
M14: $\bar A_{X,US,b}$ and $\bar A_{L,US,b}$.
&
\textbf{Proposed reference scenario:}
$\bar A_{X,US,u}N_\dagger^{\xi_u}$ and
$\bar A_{L,US,u}N_\dagger^{\nu_u}$ at a declared reference switch stock $N_\dagger$.
&
A continuity assumption at one reference state, not an estimated post-switch technology.
Uninterrupted $u$ observations cannot identify these counterfactual levels.
&
Robustness: separate proportional level shocks $m_X,m_L$ around the reference.
No empirical bounds verified. Optional $m_j\in\{0.8,1,1.2\}$ is an analyst stress grid only.
Constant bars cannot impose continuity at every possible random switch date.
\\[5pt]
\end{evidence}

\textbf{Interpretation of the post-switch scenario.}
Reference continuity pertains to factor augmentations only. Output, wages and equity values
may still jump when research allocation or expectations change.
The $(m_X,m_L)$ grid is not a literature-derived plausible interval and is not recommended
as a substitute for identifying the economic transition. Its role is to expose how conclusions
depend on an unobserved technology level. Reassess the internally calibrated financial block
for each maintained external scenario.


\clearpage
\section[Proposed candidate set: recommendations, not additional estimates]{Proposed candidate set: recommendations, not additional estimates}

All 19 included scalars appear once below. These are conditional starting choices for
research, not an operational joint calibration vector. A robustness envelope joins
measurement/specification alternatives only where stated; it is not a pooled statistical
interval. Source-specific uncertainty remains in the preceding evidence tables.
Confidence concerns \emph{economic numerical transfer}, unless algebra is separately mentioned.

\begin{recommendations}
% SCALAR: beta
M01 $\beta$
& Empirical point unassigned. Conditional endpoint bridge
$\delta_a^\Delta/(1+\delta_a^\Delta)$; HKT $0.5$ for comparison.
& Exponential sources and assumed $\Delta=20$--$40$: $0.1465$--$0.3070$.
Separate present-bias continuation scenario extends to $0.4454$.
No V6 confidence interval.
& Low; exact algebra after life-stage assumptions.
& \src{GP02} \src{LLMRT26} \src{HKT25}\\[5pt]
% SCALAR: gamma
M02 $\gamma$
& $5$ finance-oriented; $2$ explicit lifecycle alternative.
& Core $\{0.5,2,5,6.3,10\}$; separate recursive-utility stress
$\{17,20,60\}$. Evidence-motivated scenario grid; retain wide source intervals.
& Medium concept; low number.
& \src{VA03} \src{KSS08} \src{CFL2013} \src{EHI25}\\[5pt]
% SCALAR: pi
M07 $\pi$
& Conditional productivity-duration scenario $0.990^{4\Delta}$;
no direct V6 estimate.
& Mean spells $D=10,25,50$ years:
$[1-1/(4D)]^{4\Delta}$.
Only 25-year center source-derived; disaster events are separate alternatives.
& Low event mapping; exact survival conversion.
& \src{KR07} \src{Barro06} \src{BU08}\\[5pt]
% SCALAR: vartheta_US
M10 $\vartheta_{US}$
& $0.621$ (COGS-based, 2016); competing cost-definition candidate $0.870$.
& $0.621$--$0.909$ measurement-specification envelope;
$0.543$ global-database US alternative.
Not a sampling confidence range.
& Low--medium; reciprocal exact.
& \src{DLEU20} \src{Traina18} \src{DE21}\\[5pt]
% SCALAR: vartheta_W
M10 $\vartheta_W$
& Full-RoW point unassigned. Europe $0.613$ and Asia $0.690$ are geographic anchors.
& Regional proxies $0.613$--$0.725$;
coverage and markup-method sensitivity required before aggregation.
This is not a RoW interval.
& Low; aggregate not matched.
& \src{DE21}\\[5pt]
% SCALAR: rho_US
M12 $\rho_{US}$
& Skill proxy $0.667$; retain meta candidates $0.16$ (US), $0.27$ (pooled).
& Skill branch $0.16$--$0.91$, stress $0.05$.
Separate capital--labor analogy: $2$, range $1.43$--$3.33$.
No pooling across input pairs.
& Qualified skill proxy; low exact V6.
& \src{KM92} \src{CP05} \src{HILZ24} \src{OR21} \src{GHIK22}\\[5pt]
% SCALAR: rho_W
M12 $\rho_W$
& $0.667$ only under explicit US-skill pooling;
pooled/developing meta alternatives $0.27/0.47$.
& Same broad skill branch conditional on coverage.
Complementarity branch $2$ with $1.43$--$3.33$ is an analogy, not a RoW estimate.
& Low geographic and input mapping.
& \src{CP05} \src{CC06} \src{HILZ24} \src{OR21} \src{GHIK22}\\[5pt]
% SCALAR: a_US
M09 $a_{US}$
& $(G_{N,US}-1)/(s_{R,US}H_{US})$ after matched measurement, $s_{R,US}>0$.
& No portable numerical range. Propagate joint growth/research-share bounds.
HKT $aH=0.2$ is illustration only.
& Low empirical; conditional identity exact.
& \src{HKT25} \src{BJRW20} \src{BP07}\\[5pt]
% SCALAR: a_W
M09 $a_W$
& $(G_{N,W}-1)/(s_{R,W}H_W)$ for declared RoW, $s_{R,W}>0$.
& Country coverage, knowledge proxy and research-share scenarios;
no evidence for imposing the US level.
& Low; RoW aggregation unresolved.
& \src{BP07} \src{BJRW20}\\[5pt]
\end{recommendations}

\textbf{Why these numeric candidates.}
The choice $\gamma=5$ gives a provisional finance-oriented point while preserving both
lower lifecycle and much higher recursive-utility evidence. The skill candidate
$\rho=2/3$ prioritizes an original instrumental-variable study; newer bias-corrected
meta-estimates receive explicit alternatives rather than being averaged away.
The two markup candidates reflect incompatible cost definitions.
These judgments do not establish identification or override V6's restrictions.

\clearpage
\subsection{Remaining Group B scalars: normalization, data inversion and scenarios}

Here $B_X,B_L$ are the explicit V6 inversions in equations~(\ref{eq:bx})--(\ref{eq:bl});
$N_\dagger$ is a declared reference switch stock.
An unassigned empirical value is a substantive finding, not permission to replace it
with a convenient number. Suggested normalizations and stress grids are marked accordingly.

\begin{recommendations}
% SCALAR: alpha_US
M11 $\alpha_{US}$
& $0.5$ as a normalization, jointly recovering US scales, for $\rho_{US}\ne1$.
& Reparameterization checks preserving both CES composites;
no independently meaningful raw-weight range.
At $\rho=1$, use the Cobb--Douglas share.
& High equivalence; no empirical point claim.
& \src{CL12} \src{KMW12} \src{CC06}\\[5pt]
% SCALAR: alpha_W
M11 $\alpha_W$
& $0.5$ under the same normalization logic, for $\rho_W\ne1$.
& Jointly adjust both RoW augmentations.
Near $\rho=1$, inversion sensitivity needs explicit treatment.
& High equivalence; units conditional.
& \src{CL12} \src{KMW12}\\[5pt]
% SCALAR: A_X_US_u
M13 $\bar A_{X,US,u}$
& $B_{X,US}/N_{US,0}^{\xi_u}$; numerical value awaits matched data.
& Joint output/share, $\varphi$, markup, $\rho$ and knowledge-normalization uncertainty.
No portable scale interval.
& Conditional algebra high; empirical low.
& \src{CC06} \src{Rossi22} \src{LMW10}\\[5pt]
% SCALAR: A_L_US_u
M13 $\bar A_{L,US,u}$
& $B_{L,US}/N_{US,0}^{\nu_u}$; numerical value awaits matched data.
& Joint production-data region, coupled to the preceding coefficient.
HKT's $A_LL=1$ is a units-dependent illustration.
& Conditional algebra high; empirical low.
& \src{CC06} \src{HKT25} \src{LMW10}\\[5pt]
% SCALAR: Abar_X_US
M14 $\bar A_{X,US,b}$
& $\bar A_{X,US,u}N_\dagger^{\xi_u}$: reference-state augmentation continuity scenario.
& Multiply by $m_X$; $\{0.8,1,1.2\}$ optional analyst stress grid.
No empirical bound verified.
& Low; counterfactual assumption.
& V6 equations; \src{HKT25} comparison only\\[5pt]
% SCALAR: Abar_L_US
M14 $\bar A_{L,US,b}$
& $\bar A_{L,US,u}N_\dagger^{\nu_u}$: reference-state augmentation continuity scenario.
& Multiply by $m_L$ jointly with $m_X$.
Grid is not a credible interval; cannot ensure continuity at all switch dates.
& Low; counterfactual assumption.
& V6 equations; \src{HKT25} comparison only\\[5pt]
% SCALAR: Abar_X_W
M15 $\bar A_{X,W}$
& $B_{X,W}$ for matched RoW production observations.
& Country coverage and productivity/input measurement region,
preserving relative US--RoW units.
& Low empirical; conditional inversion exact.
& \src{CC06} \src{Rossi22}\\[5pt]
% SCALAR: Abar_L_W
M15 $\bar A_{L,W}$
& $B_{L,W}$ for the same matched RoW observations.
& Joint with $\bar A_{X,W}$; no independent country-average TFP range.
& Low empirical; conditional inversion exact.
& \src{CC06} \src{Rossi22}\\[5pt]
% SCALAR: xi_u
M16 $\xi_u$
& Empirical point unassigned; HKT $0.7$ for comparison.
& Joint augmentation/knowledge-growth region.
No verified transferable numeric bounds; annual demand trends are not this elasticity.
& High HKT definition match; low empirical.
& \src{HKT25} \src{KM92} \src{BJRW20}\\[5pt]
% SCALAR: nu_u
M16 $\nu_u$
& Empirical point unassigned; HKT $0.2$ for comparison.
& Joint with $\xi_u$; preserve data-supported zero/equal/reversed growth
even outside the current theorem/solver class.
& High HKT definition match; low empirical.
& \src{HKT25} \src{KM92} \src{BP07}\\[5pt]
\end{recommendations}

\textbf{Use in subsequent calibration.}
First declare the period and labor/knowledge/output units, then choose and defend the
production-input mapping. Assign technology levels jointly with the matched production
moments, recording those moments as targeted. For each retained external specification,
re-fit Group C to the chosen international-finance moments and assess sensitivity and
identification. Data/model disagreement is evidence about the mapping or specification;
it must not be concealed by selecting only externally convenient parameters.


\clearpage
\section{References and verified versions}
Links in each entry lead to the DOI and the paper/version used for the checks. Page and table locators refer to that version where specified. Working papers, illustrative calibrations and meta-analyses retain their distinct status.
\begin{multicols}{2}\fontsize{9.2}{10.8}\selectfont\setlength{\parskip}{7pt}
\begin{minipage}{\linewidth}\hypertarget{ref:HKT25}{}\textbf{HKT25}. Tomohiro Hirano, Keiichi Kishi, Alexis Akira Toda (2025). Technological Innovation and Bursting Bubbles. \emph{arXiv:2501.08215v2, preprint}. \href{https://doi.org/10.48550/arXiv.2501.08215}{DOI}. \href{https://arxiv.org/html/2501.08215v2}{Verified source/version}. Pinned arXiv v2 and matching local original TeX; Figure 1, p. 24. Numerical illustration, not an estimate.\end{minipage}\par
\begin{minipage}{\linewidth}\hypertarget{ref:GP02}{}\textbf{GP02}. Gourinchas, Pierre-Olivier and Parker, Jonathan A. (2002). Consumption Over the Life Cycle. \emph{Econometrica 70(1), 47-89}. \href{https://doi.org/10.1111/1468-0262.00269}{DOI}. \href{https://jggarita.github.io/MacroIntermedia/materials/Readings/GourinchasParker2002.pdf}{Verified source/version}. Table III p.71; Table V p.79\end{minipage}\par
\begin{minipage}{\linewidth}\hypertarget{ref:LLMRT26}{}\textbf{LLMRT26}. Laibson, David and Lee, Sean Chanwook and Maxted, Peter and Repetto, Andrea and Tobacman, Jeremy (2026). Estimating Discount Functions with Consumption Choices over the Lifecycle. \emph{Review of Financial Studies 39(6), 1580-1610; online 2024, issue 2026}. \href{https://doi.org/10.1093/rfs/hhae035}{DOI}. \href{https://academic.oup.com/rfs/article/39/6/1580/7703274?rss=1}{Verified source/version}. Table 3; footnotes 25 and 39\end{minipage}\par
\begin{minipage}{\linewidth}\hypertarget{ref:BJKS97}{}\textbf{BJKS97}. Barsky, Robert and Juster, F. Thomas and Kimball, Miles and Shapiro, M. (1997). Preference Parameters and Behavioral Heterogeneity: An Experimental Approach in the Health and Retirement Study. \emph{Quarterly Journal of Economics 112(2), 537-579}. \href{https://doi.org/10.1162/003355397555280}{DOI}. \href{https://websites.umich.edu/~mkimball/keio/5.\%20papers/qje-bjks-risk-aversion.pdf}{Verified source/version}. Table XI p.563\end{minipage}\par
\begin{minipage}{\linewidth}\hypertarget{ref:KSS08}{}\textbf{KSS08}. Kimball, Miles and Sahm, Claudia and Shapiro, Matthew (2008). Imputing Risk Tolerance From Survey Responses. \emph{Journal of the American Statistical Association 103(483), 1028-1038}. \href{https://doi.org/10.1198/016214508000000139}{DOI}. \href{https://pmc.ncbi.nlm.nih.gov/articles/PMC2856097/}{Verified source/version}. Equations 1--2; Table 4\end{minipage}\par
\begin{minipage}{\linewidth}\hypertarget{ref:HL02}{}\textbf{HL02}. Holt, Charles A. and Laury, Susan K. (2002). Risk Aversion and Incentive Effects. \emph{American Economic Review 92(5), 1644-1655}. \href{https://doi.org/10.1257/000282802762024700}{DOI}. \href{https://lorko.sk/wp-content/uploads/2023/10/HoltLaury2002.pdf}{Verified source/version}. Table 3 and discussion p.1649\end{minipage}\par
\begin{minipage}{\linewidth}\hypertarget{ref:VA03}{}\textbf{VA03}. Vissing-Jørgensen, Annette and Attanasio, Orazio (2003). Stock-Market Participation, Intertemporal Substitution, and Risk-Aversion. \emph{American Economic Review Papers and Proceedings 93(2), 383-391}. \href{https://doi.org/10.1257/000282803321947399}{DOI}. \href{https://www.researchgate.net/publication/4981014_Stock-Market_Participation_Intertemporal_Substitution_and_Risk-Aversion}{Verified source/version}. Numeric evidence: author-uploaded extended manuscript dated 10 January 2003, Table 2A, case 4, p. 14; journal article is shorter.\end{minipage}\par
\begin{minipage}{\linewidth}\hypertarget{ref:CFL2013}{}\textbf{CFL2013}. Xiaohong Chen, Jack Favilukis, Sydney C. Ludvigson (2013). An estimation of economic models with recursive preferences. \emph{Quantitative Economics 4(1), 39-83}. \href{https://doi.org/10.3982/QE97}{DOI}. \href{https://www.econstor.eu/bitstream/10419/150343/1/v4-i1-p039-083-QE97.pdf}{Verified source/version}. Equations (1)-(2), p. 44; Table 2, p. 62; methods and interpretation, pp. 56-58, 64-65.\end{minipage}\par
\begin{minipage}{\linewidth}\hypertarget{ref:EHI25}{}\textbf{EHI25}. Elminejad, Ali and Havranek, Tomas and Irsova, Zuzana (2025). Relative Risk Aversion: A Meta-Analysis. \emph{Journal of Economic Surveys 39(5), 2315-2333}. \href{https://doi.org/10.1111/joes.12689}{DOI}. \href{https://onlinelibrary.wiley.com/doi/full/10.1111/joes.12689}{Verified source/version}. Sections 1/5 and Table 6\end{minipage}\par
\begin{minipage}{\linewidth}\hypertarget{ref:DLEU20}{}\textbf{DLEU20}. Jan De Loecker, Jan Eeckhout and Gabriel Unger (2020). The Rise of Market Power and the Macroeconomic Implications. \emph{Quarterly Journal of Economics 135(2), 561--644}. \href{https://doi.org/10.1093/qje/qjz041}{DOI}. \href{https://www.janeeckhout.com/wp-content/uploads/RMP.pdf}{Verified source/version}. Author manuscript, 15 November 2019: equation (7), section 3.1 and Figure 1, pp. 9-10.\end{minipage}\par
\begin{minipage}{\linewidth}\hypertarget{ref:Traina18}{}\textbf{Traina18}. James Traina (2018). Is Aggregate Market Power Increasing? Production Trends Using Financial Statements. \emph{Working paper, February2018}. \href{https://doi.org/10.2139/ssrn.3120849}{DOI}. \href{https://doi.org/10.2139/ssrn.3120849}{Verified source/version}. Longer Time Horizons subsection and Figure 5; numerical passage verified in original full text via Scite.\end{minipage}\par
\begin{minipage}{\linewidth}\hypertarget{ref:DE21}{}\textbf{DE21}. Jan De Loecker and Jan Eeckhout (2021). Global Market Power. \emph{Working paper, 10 February 2021 version; NBER Working Paper 24768 first issued 2018}. \href{https://doi.org/10.3386/w24768}{DOI}. \href{https://crei.cat/wp-content/uploads/2021/04/GLOBAL.pdf}{Verified source/version}. Data: section 2; regional estimates: Table 1, manuscript p. 13.\end{minipage}\par
\begin{minipage}{\linewidth}\hypertarget{ref:BHKZ21}{}\textbf{BHKZ21}. Steve Bond, Arshia Hashemi, Greg Kaplan and Piotr Zoch (2021). Some Unpleasant Markup Arithmetic: Production Function Elasticities and Their Estimation from Production Data. \emph{Journal of Monetary Economics 121, 1-14}. \href{https://doi.org/10.1016/j.jmoneco.2021.05.004}{DOI}. \href{https://arshiahashemi.com/publication/jme-2021-bhkz/JME-2021-BHKZ.pdf}{Verified source/version}. Abstract verified; published-PDF direct fetch timed out\end{minipage}\par
\begin{minipage}{\linewidth}\hypertarget{ref:DL21}{}\textbf{DL21}. Jan De Loecker (2021). Comment on (Un)pleasant ... By Bond Et Al (2020). \emph{Journal of Monetary Economics 121, 15-18}. \href{https://doi.org/10.1016/j.jmoneco.2021.04.009}{DOI}. \href{https://doi.org/10.1016/j.jmoneco.2021.04.009}{Verified source/version}. Introduction, Scite fulltext\end{minipage}\par
\begin{minipage}{\linewidth}\hypertarget{ref:KM92}{}\textbf{KM92}. Lawrence F. Katz and Kevin M. Murphy (1992). Changes in Relative Wages, 1963--1987: Supply and Demand Factors. \emph{Quarterly Journal of Economics 107(1), 35--78}. \href{https://doi.org/10.2307/2118323}{DOI}. \href{https://lkatz.scholars.harvard.edu/sites/g/files/omnuum5961/files/lkatz/files/changes_in_relative_wages_1963-1987_supply_and_demand_factors.pdf}{Verified source/version}. Equation (19), p.69; elasticity discussion p.72\end{minipage}\par
\begin{minipage}{\linewidth}\hypertarget{ref:CP05}{}\textbf{CP05}. Antonio Ciccone and Giovanni Peri (2005). Long-Run Substitutability Between More and Less Educated Workers: Evidence from U.S. States, 1950--1990. \emph{Review of Economics and Statistics 87(4), 652--663}. \href{https://doi.org/10.1162/003465305775098233}{DOI}. \href{https://giovanniperi.ucdavis.edu/uploads/5/6/8/2/56826033/ciccone_peri_long_run_substitutability_\%40005.pdf}{Verified source/version}. Table 4, p.656, panel B(iv), column 1; full panels A--C for estimator sensitivity\end{minipage}\par
\begin{minipage}{\linewidth}\hypertarget{ref:KORV00}{}\textbf{KORV00}. Per Krusell, Lee E. Ohanian, Jose-Victor Rios-Rull and Giovanni L. Violante (2000). Capital-Skill Complementarity and Inequality: A Macroeconomic Analysis. \emph{Econometrica 68(5), 1029--1053}. \href{https://doi.org/10.1111/1468-0262.00150}{DOI}. \href{https://www.jonathanheathcote.com/KORV.pdf}{Verified source/version}. Table II, p. 1041; alternative skill threshold: footnote 16, p. 1044.\end{minipage}\par
\begin{minipage}{\linewidth}\hypertarget{ref:HILZ24}{}\textbf{HILZ24}. Tomas Havranek, Zuzana Irsova, Lubica Laslopova and Olesia Zeynalova (2024). Publication and Attenuation Biases in Measuring Skill Substitution. \emph{Review of Economics and Statistics 106(5), 1187-1200; online 2022}. \href{https://doi.org/10.1162/rest_a_01227}{DOI}. \href{https://meta-analysis.cz/skill/skill2.pdf}{Verified source/version}. Final published Table 5, p. 1197; supersedes 2020/2021 low-elasticity working paper.\end{minipage}\par
\begin{minipage}{\linewidth}\hypertarget{ref:CC06}{}\textbf{CC06}. Francesco Caselli and Wilbur John Coleman II (2006). The World Technology Frontier. \emph{American Economic Review 96(3), 499--522}. \href{https://doi.org/10.1257/aer.96.3.499}{DOI}. \href{https://personal.lse.ac.uk/casellif/papers/wtf.pdf}{Verified source/version}. Equations (2)-(3), pp. 499-502; section I.B and Table 2, pp. 504-505.\end{minipage}\par
\begin{minipage}{\linewidth}\hypertarget{ref:OR21}{}\textbf{OR21}. Ezra Oberfield and Devesh Raval (2021). Micro Data and Macro Technology. \emph{Econometrica 89(2), 703-732}. \href{https://doi.org/10.3982/ECTA12807}{DOI}. \href{https://ezraoberfield.github.io/CESAggregation.pdf}{Verified source/version}. Abstract, p. 703; introduction, p. 705; Figure 4b, p. 719.\end{minipage}\par
\begin{minipage}{\linewidth}\hypertarget{ref:GHIK22}{}\textbf{GHIK22}. Sebastian Gechert, Tomas Havranek, Zuzana Irsova and Dominika Kolcunova (2022). Measuring Capital-Labor Substitution: The Importance of Method Choices and Publication Bias. \emph{Review of Economic Dynamics 45, 55-82}. \href{https://doi.org/10.1016/j.red.2021.05.003}{DOI}. \href{https://meta-analysis.cz/sigma/sigma2.pdf}{Verified source/version}. Table 9, p. 77; method: section 5.4; conclusion, p. 78.\end{minipage}\par
\begin{minipage}{\linewidth}\hypertarget{ref:KR07}{}\textbf{KR07}. James A. Kahn and Robert W. Rich (2007). Tracking the new economy: Using growth theory to detect changes in trend productivity. \emph{Journal of Monetary Economics 54(6), 1670-1701}. \href{https://doi.org/10.1016/j.jmoneco.2006.07.008}{DOI}. \href{https://pages.stern.nyu.edu/~jkahn/Papers/kahnrich_jme.pdf}{Verified source/version}. Equations (7)-(9), p. 1677; Table 2, p. 1683; sample/durations, p. 1682; identification, p. 1687.\end{minipage}\par
\begin{minipage}{\linewidth}\hypertarget{ref:Barro06}{}\textbf{Barro06}. Robert J. Barro (2006). Rare Disasters and Asset Markets in the Twentieth Century. \emph{Quarterly Journal of Economics 121(3), 823-866}. \href{https://doi.org/10.1162/qjec.121.3.823}{DOI}. \href{https://dash.harvard.edu/bitstreams/7312037c-67b1-6bd4-e053-0100007fdf3b/download}{Verified source/version}. p.825 disaster process; p.831 observed annual frequency; Table V, p.846\end{minipage}\par
\begin{minipage}{\linewidth}\hypertarget{ref:BU08}{}\textbf{BU08}. Robert J. Barro, Jose F. Ursua (2008). Macroeconomic Crises since 1870. \emph{Brookings Papers on Economic Activity, Spring, 255-350 including discussion}. \href{https://doi.org/10.1353/eca.0.0000}{DOI}. \href{https://www.brookings.edu/wp-content/uploads/2008/03/2008a_bpea_barro.pdf}{Verified source/version}. pp.284-286 hazard definition and sample; Table 10, p.296; Table 11, p.297; footnote24, p.286\end{minipage}\par
\begin{minipage}{\linewidth}\hypertarget{ref:BJRW20}{}\textbf{BJRW20}. Nicholas Bloom, Charles I. Jones, John Van Reenen, Michael Webb (2020). Are Ideas Getting Harder to Find?. \emph{American Economic Review 110(4), 1104-1144}. \href{https://doi.org/10.1257/aer.20180338}{DOI}. \href{https://web.stanford.edu/~chadj/IdeaPF.pdf}{Verified source/version}. Equations (1)-(3), pp. 1108-1110; Table 7 and equation (17), p. 1134; Table A1, p. 1140.\end{minipage}\par
\begin{minipage}{\linewidth}\hypertarget{ref:BP07}{}\textbf{BP07}. Laura Bottazzi and Giovanni Peri (2007). The International Dynamics of R\&D and Innovation in the Long Run and in the Short Run. \emph{Economic Journal 117(518), 486-511}. \href{https://doi.org/10.1111/j.1468-0297.2007.02027.x}{DOI}. \href{https://giovanniperi.ucdavis.edu/uploads/5/6/8/2/56826033/bottazzi_et_al-2007-the_economic_journal.pdf}{Verified source/version}. Equations (5)-(9), pp. 492-494; sample, p. 495; Table 4, p. 501; interpretation, p. 502.\end{minipage}\par
\begin{minipage}{\linewidth}\hypertarget{ref:Jones95}{}\textbf{Jones95}. Charles I. Jones (1995). R \& D-Based Models of Economic Growth. \emph{Journal of Political Economy 103(4), 759-784}. \href{https://doi.org/10.1086/262002}{DOI}. \href{https://www.journals.uchicago.edu/doi/10.1086/262002}{Verified source/version}. publisher abstract; qualitative specification critique only\end{minipage}\par
\begin{minipage}{\linewidth}\hypertarget{ref:Rossi22}{}\textbf{Rossi22}. Federico Rossi (2022). The Relative Efficiency of Skilled Labor across Countries: Measurement and Interpretation. \emph{American Economic Review 112(1), 235-266}. \href{https://doi.org/10.1257/aer.20191852}{DOI}. \href{https://wrap.warwick.ac.uk/id/eprint/159013/1/WRAP-Relative-efficiency-skilled-labor-across-countries-measurement-interpretation-2021.pdf}{Verified source/version}. July 2021 accepted manuscript, equations (4)-(8), pp. 7-8; Table I, p. 31 (PDF page 32). Publication metadata checked against AEA.\end{minipage}\par
\begin{minipage}{\linewidth}\hypertarget{ref:CL12}{}\textbf{CL12}. Cristiano Cantore and Paul Levine (2012). Getting Normalization Right: Dealing with Dimensional Constants in Macroeconomics. \emph{Journal of Economic Dynamics and Control 36(12), 1931-1949}. \href{https://doi.org/10.1016/j.jedc.2012.05.009}{DOI}. \href{https://archives.dynare.org/wp-repo/dynarewp009.pdf}{Verified source/version}. Detailed method: Dynare Working Paper 9 (July 2011), sections 2.1-2.5.\end{minipage}\par
\begin{minipage}{\linewidth}\hypertarget{ref:KMW12}{}\textbf{KMW12}. Rainer Klump, Peter McAdam and Alpo Willman (2012). The Normalized CES Production Function: Theory and Empirics. \emph{Journal of Economic Surveys 26(5), 769-799}. \href{https://doi.org/10.1111/j.1467-6419.2012.00730.x}{DOI}. \href{https://www.ecb.europa.eu/pub/pdf/scpwps/ecbwp1294.pdf}{Verified source/version}. Detailed method: ECB Working Paper 1294 (2011), equations (10), (15)-(16).\end{minipage}\par
\begin{minipage}{\linewidth}\hypertarget{ref:LMW10}{}\textbf{LMW10}. Miguel A. Leon-Ledesma, Peter McAdam and Alpo Willman (2010). Identifying the Elasticity of Substitution with Biased Technical Change. \emph{American Economic Review 100(4), 1330-1357}. \href{https://doi.org/10.1257/aer.100.4.1330}{DOI}. \href{https://www.ecb.europa.eu/pub/pdf/scpwps/ecbwp1001.pdf}{Verified source/version}. Method verified in ECB Working Paper 1001 (January 2009); published metadata checked at AEA.\end{minipage}\par
\end{multicols}
\end{document}
